%------------------------------------------------------------------------------%
\begin{question}
  Considerando a equação
  \[
    8(x + 1) + (y - 3)^2 = 0
  \]
  \vspace{-\baselineskip}
  \begin{enumerate}
    \item identifique a curva
    \item determine seus principais valores e pontos
    \item esboce a curva
  \end{enumerate}
\end{question}

%------------------------------------------------------------------------------%
\begin{resolution}{Identificação}
\begin{columns}[t]
\column{0.4\textwidth}
  \onslide<+->{A equação
  \[
    (y - 3)^2 = - 8(x + 1)
  \]}%
  \onslide<+->{está na forma de uma parábola transladada
  \[
    (y - k)^2 = - 8(x + h)
  \]}
  \onslide<+->{Identificamos o vértice}
  \[
    \onslide<+->{C = ( h,  k)}
    \onslide<+->{  = (-1, -3)}
  \]
\column{0.4\textwidth}
  \[
    \onslide<+->{4p = -8}
    \qquad
    \onslide<+->{p = -2}
  \]

  \bigskip
  \onslide<+->{O termo elevado ao quadrado é $(y-k)^2$,}
  \onslide<+->{o eixo de simetria é horizontal}

  \bigskip
  \onslide<+->{Como $p<0$ a parábola tem abertura para a esquerda}
\end{columns}
\end{resolution}

%------------------------------------------------------------------------------%
\begin{resolution}{Esboço}
  \centering
  \includegraphics{I-07-Identificacao-ex-9}
\end{resolution}

%------------------------------------------------------------------------------%
